% Options for packages loaded elsewhere
\PassOptionsToPackage{unicode}{hyperref}
\PassOptionsToPackage{hyphens}{url}
\documentclass[
  spanish,
  11pt,
]{article}
\usepackage{xcolor}
\usepackage[margin=2.5cm]{geometry}
\usepackage{amsmath,amssymb}
\setcounter{secnumdepth}{5}
\usepackage{iftex}
\ifPDFTeX
  \usepackage[T1]{fontenc}
  \usepackage[utf8]{inputenc}
  \usepackage{textcomp} % provide euro and other symbols
\else % if luatex or xetex
  \usepackage{unicode-math} % this also loads fontspec
  \defaultfontfeatures{Scale=MatchLowercase}
  \defaultfontfeatures[\rmfamily]{Ligatures=TeX,Scale=1}
\fi
\usepackage{lmodern}
\ifPDFTeX\else
  % xetex/luatex font selection
\fi
% Use upquote if available, for straight quotes in verbatim environments
\IfFileExists{upquote.sty}{\usepackage{upquote}}{}
\IfFileExists{microtype.sty}{% use microtype if available
  \usepackage[]{microtype}
  \UseMicrotypeSet[protrusion]{basicmath} % disable protrusion for tt fonts
}{}
\makeatletter
\@ifundefined{KOMAClassName}{% if non-KOMA class
  \IfFileExists{parskip.sty}{%
    \usepackage{parskip}
  }{% else
    \setlength{\parindent}{0pt}
    \setlength{\parskip}{6pt plus 2pt minus 1pt}}
}{% if KOMA class
  \KOMAoptions{parskip=half}}
\makeatother
\usepackage{graphicx}
\makeatletter
\newsavebox\pandoc@box
\newcommand*\pandocbounded[1]{% scales image to fit in text height/width
  \sbox\pandoc@box{#1}%
  \Gscale@div\@tempa{\textheight}{\dimexpr\ht\pandoc@box+\dp\pandoc@box\relax}%
  \Gscale@div\@tempb{\linewidth}{\wd\pandoc@box}%
  \ifdim\@tempb\p@<\@tempa\p@\let\@tempa\@tempb\fi% select the smaller of both
  \ifdim\@tempa\p@<\p@\scalebox{\@tempa}{\usebox\pandoc@box}%
  \else\usebox{\pandoc@box}%
  \fi%
}
% Set default figure placement to htbp
\def\fps@figure{htbp}
\makeatother
% definitions for citeproc citations
\NewDocumentCommand\citeproctext{}{}
\NewDocumentCommand\citeproc{mm}{%
  \begingroup\def\citeproctext{#2}\cite{#1}\endgroup}
\makeatletter
 % allow citations to break across lines
 \let\@cite@ofmt\@firstofone
 % avoid brackets around text for \cite:
 \def\@biblabel#1{}
 \def\@cite#1#2{{#1\if@tempswa , #2\fi}}
\makeatother
\newlength{\cslhangindent}
\setlength{\cslhangindent}{1.5em}
\newlength{\csllabelwidth}
\setlength{\csllabelwidth}{3em}
\newenvironment{CSLReferences}[2] % #1 hanging-indent, #2 entry-spacing
 {\begin{list}{}{%
  \setlength{\itemindent}{0pt}
  \setlength{\leftmargin}{0pt}
  \setlength{\parsep}{0pt}
  % turn on hanging indent if param 1 is 1
  \ifodd #1
   \setlength{\leftmargin}{\cslhangindent}
   \setlength{\itemindent}{-1\cslhangindent}
  \fi
  % set entry spacing
  \setlength{\itemsep}{#2\baselineskip}}}
 {\end{list}}
\usepackage{calc}
\newcommand{\CSLBlock}[1]{\hfill\break\parbox[t]{\linewidth}{\strut\ignorespaces#1\strut}}
\newcommand{\CSLLeftMargin}[1]{\parbox[t]{\csllabelwidth}{\strut#1\strut}}
\newcommand{\CSLRightInline}[1]{\parbox[t]{\linewidth - \csllabelwidth}{\strut#1\strut}}
\newcommand{\CSLIndent}[1]{\hspace{\cslhangindent}#1}
\ifLuaTeX
\usepackage[bidi=basic,shorthands=off]{babel}
\else
\usepackage[bidi=default,shorthands=off]{babel}
\fi
\ifLuaTeX
  \usepackage{selnolig} % disable illegal ligatures
\fi
\setlength{\emergencystretch}{3em} % prevent overfull lines
\providecommand{\tightlist}{%
  \setlength{\itemsep}{0pt}\setlength{\parskip}{0pt}}
\usepackage{graphicx}
\usepackage{float}
\usepackage[default]{cantarell}
\usepackage{setspace}
\onehalfspacing
\usepackage{svg}
\usepackage{bookmark}
\IfFileExists{xurl.sty}{\usepackage{xurl}}{} % add URL line breaks if available
\urlstyle{same}
% fallback for those not using the hyperref driver hyperxmp:
\makeatletter
\@ifundefined{xmpquote}{\newcommand{\xmpquote}[1]{#1}}{}
\makeatother
\hypersetup{
  pdflang={es-AR},
  hidelinks,
  pdfcreator={LaTeX via pandoc}}

\author{}
\date{\vspace{-2.5em}}

\begin{document}

\begin{titlepage}
    \centering
    \includegraphics[width=.3\textwidth]{../../../img/udemm-logo.png}
    
    \vspace{0.1cm}
    {\Large Facultad de Ingeniería} \\
    \vspace{0.2cm}
    {\large Ingeniería en Sistemas} \\
    
    \vspace{2.5cm}
    \rule{\linewidth}{0.4mm} \\[0.4cm]
    {\huge \bfseries Trabajo Práctico} \\[0.3cm]
    {\Large Arquitectura Cliente-Servidor} \\[0.2cm]
    \rule{\linewidth}{0.4mm}
    
    \vfill
    
    \begin{tabular}{r l}
        \textbf{Asignatura:} & Diseño de Sistemas II \\
        \textbf{Docente:}    & Mg. Daniel Omar García \\
        \textbf{Alumno:}     & Daniel Ise \\
        \textbf{Legajo:}     & 28547 \\
        \textbf{Fecha:}      & Noviembre, 2026
    \end{tabular}
\end{titlepage}

\pagenumbering{roman}
\thispagestyle{empty}
\tableofcontents
\newpage

\pagenumbering{arabic}
\setcounter{page}{1}

\section{Introducción}\label{introducciuxf3n}

La arquitectura Cliente-Servidor es un modelo de diseño de
\emph{software} basado en la división del sistema en dos componentes
principales: el \emph{cliente} y el \emph{servidor}.\footnote{Daniel
  Omar García, \emph{Arquitectura Cliente/Servidor}, Apunte de cátedra,
  Diseño de Sistemas II, Universidad de la Marina Mercante, 2026.} Su
funcionamiento sigue el \emph{ciclo de petición y respuesta}:

\begin{itemize}
\tightlist
\item
  El \textbf{cliente} envía una \emph{petición} al servidor solicitando
  un recurso o servicio.
\item
  El \textbf{servidor} \emph{procesa} la solicitud y elabora una
  \emph{respuesta}.
\item
  El \textbf{servidor} \emph{envía} esa respuesta al cliente.
\item
  El \textbf{cliente} \emph{recibe} la respuesta y la presenta o procesa
  localmente.
\end{itemize}

Toda la comunicación se rige por un \emph{protocolo} que especifica las
\emph{reglas} de que cliente y servidor deben seguir. Por ejemplo, el
cliente puede ser el navegador web \texttt{Chromium}, que envía una
petición a un servidor web \texttt{Apache}. El servidor recibe la
solicitud, ejecuta la lógica correspondiente en \texttt{PHP} y genera
una respuesta en formato \texttt{HTML}, que envía de regreso al cliente
para que este la renderice en pantalla. Este intercambio se realiza
siguiendo las reglas de los protocolos \texttt{TCP/IP} y \texttt{HTTP}.

Adoptar una arquitectura Cliente-Servidor ofrece algunas
\textbf{ventajas}.\footnote{García, \emph{Arquitectura
  Cliente/Servidor}.} En primer lugar, la \textbf{centralización}, dado
que los recursos y datos relevantes del sistema residen en el servidor.
Segundo, facilita la \textbf{administración} de la información desde un
único punto, aspecto clave para aplicaciones que deben cumplir, por
ejemplo, normativas de protección de datos. Tercero, permite
\textbf{escalar} la capacidad de procesamiento para atender a más
clientes ante incrementos de tráfico. Por último, posibilita un mejor
\textbf{control de la seguridad} mediante políticas centralizadas en el
servidor.

Este estilo arquitectónico trae también sus
\textbf{desventajas}.\footnote{García, \emph{Arquitectura
  Cliente/Servidor}.} Una aplicación basada en este modelo
\textbf{depende críticamente del servidor}: un error o caída del mismo
puede dejar a todos los usuarios sin acceso al servicio. Asimismo,
desplegar y mantener una infraestructura propia suele demandar una
\textbf{inversión significativa} en \emph{hardware} y seguridad.
Finalmente, el sistema puede sufrir \textbf{saturación} frente a picos
de usuarios concurrentes, si no fue diseñado adecuadamente para escalar.

El presente trabajo práctico analiza cuatro casos exitosos de
implementación de esta arquitectura:

\begin{enumerate}
\def\labelenumi{\arabic{enumi}.}
\tightlist
\item
  \textbf{Netflix}, plataforma de \emph{streaming} de video.
\item
  \textbf{Spotify}, servicio de \emph{streaming} de audio (en
  particular, música y \emph{podcasts}).
\item
  \textbf{WhatsApp}, aplicación de mensajería instantánea.
\item
  \textbf{Gmail}, servicio pionero de correo electrónico web.
\end{enumerate}

Se eligieron estas aplicaciones por ser referentes en la resolución de
los principales desafíos del modelo Cliente-Servidor: todas gestionan
grandes volúmenes de datos, responden a picos de demanda y garantizan la
seguridad y privacidad mediante mecanismos robustos de autenticación y
autorización. Además, operan en varias plataformas, computadoras de
escritorio, teléfonos móviles, portátiles y televisores, exigiendo altos
estándares de rendimiento, escalabilidad y experiencia de usuario.

En lo que respecta a los casos de implementaciones problemáticas, se
eligieron:

\begin{enumerate}
\def\labelenumi{\arabic{enumi}.}
\tightlist
\item
  \textbf{SIU-Guaraní}, sistema de autogestión académica utilizado en
  universidades argentinas.
\item
  \textbf{SimCity (2013)}, videojuego cuyo lanzamiento exigía conexión
  permanente a internet para jugar.
\item
  \textbf{Bolsas y Mercados Argentinos (BYMA)}, infraestructura bursátil
  argentina que ha experimentado incidentes en jornadas de alta
  operatoria.
\end{enumerate}

Estos casos fueron elegidos porque permiten examinar fallas recurrentes
del modelo frente a picos de concurrencia y dependencia crítica de la
red, evidenciando riesgos como cuellos de botella, sobrecarga del
servidor e indisponibilidad del servicio.\footnote{García,
  \emph{Arquitectura Cliente/Servidor}.}

\section{Casos exitosos}\label{casos-exitosos}

\subsection{Netflix}\label{netflix}

\textbf{Netflix} es un proveedor de servicios de \emph{streaming} de
video, enfocado en series y películas. Enfrenta grandes desafíos de
ingeniería en lo que respecta a rendimiento: debe recibir, procesar y
enviar grandes volúmenes de datos a cientos de miles de clientes de todo
el mundo, que se conectan desde dispositivos muy diferentes en
capacidades de cómputo. El Servidor debe de proveer autenticación,
catálogo, perfiles y transmisión de video en tiempo real, por lo que
supone un caso de análisis particularmente interesante.

El UML es un diagrama de secuencia representativo es un usuario
conectándose al servicio con sus credenciales, eligiendo del catálogo
para el que tiene autorización, y recibiendo el \emph{streaming} de
video.

\includesvg{uml/netflix.svg}

La arquitectura Cliente-Servidor es fundamental para un servicio de
streaming, que debe enviar grandes volúmenes de datos a muchos usuarios
en todo el mundo.

\subsection{Spotify}\label{spotify}

El segundo caso a considerar es el de \textbf{Spotify}. Se trata también
de un servicio de \emph{streaming}, en este caso de audio, siendo quizá
el servicio más exitoso en su rubro. Sus servidores deben de proveer a
los clientes, que utilizan todo tipo de dispositivos, de autenticación,
recomendaciones, acceso a catálogo y transmisión de audio. Se trata de
otro caso de una empresa de gran éxito, con una amplia variedad de
clientes (en términos de dispositivos diferenciados), por lo que su
análisis reviste de particular importancia.

\subsection{Whatsapp}\label{whatsapp}

El tercer caso de éxito a considerar es el de \textbf{Whatsapp}, una
aplicación de mensajería instantánea, orientada principal aunque no
exclusivamente a la mensajería móvil. El servidor debe de gestionar
colas de mensajes, notificaciones, autenticación y encriptación,
tratando de minimizar, en los casos en que corresponda, el consumo de
batería en sus clientes. La gran difusión de esta aplicación, así como
el desafío de ingeniería que supone el manejo del enorme flujo de datos
diarios, motiva a incluirla en este trabajo práctico.

\subsection{Gmail}\label{gmail}

Por último, el cuarto caso es el de \textbf{Gmail}, aplicación pionera
en la provisión de \emph{email} en navegador, que utiliza tecnología
AJAX del lado del cliente, para permitir el uso del correo electrónico
de forma eficiente sin requerir de la instalación de aplicaciones de
escritorio, como era el estándar del rubro en el momento de su
aparición. Su desarrollo y evolución como aplicación que implementa la
arquitectura de Cliente-Servidor parece casi obvia, por la innovación
que representó en su momento y por ser pionera en las aplicaciones web
asíncronas, modelo para las que vendrían luego.

\section{Casos problemáticos}\label{casos-problemuxe1ticos}

\subsection{SIU Guaraní}\label{siu-guaranuxed}

\subsection{SimCity}\label{simcity}

\subsection{Bolsas y Mercados Argentinos
(BYMA)}\label{bolsas-y-mercados-argentinos-byma}

\section{Conclusiones}\label{conclusiones}

\section*{Referencias
bibliográficas}\label{referencias-bibliogruxe1ficas}
\addcontentsline{toc}{section}{Referencias bibliográficas}

\protect\phantomsection\label{refs}
\begin{CSLReferences}{1}{1}
\bibitem[\citeproctext]{ref-garcia2026}
García, Daniel Omar. \emph{Arquitectura Cliente/Servidor}. Apunte de
cátedra, Diseño de Sistemas II, Universidad de la Marina Mercante, 2026.

\end{CSLReferences}

\end{document}
